Summing up, we have stated:
\begin{enumerate}
    \item A function is not a mapping operation, it is basically a relation between pairs composed by values coming from a domain and a codomain.
    \item Any function is usually described by: $$f: A \rightarrow B$$ where:
    \begin{itemize}
        \renewcommand{\labelitemi}{-}
        \item $A \rightarrow$ domain
        \item $B \rightarrow$ codomain
    \end{itemize}
    \item Sometimes $A$ and $B$ are countable not-string sets, therefore we must encode them following some strategies.
    \item The encoding of any element $x$ is indicated as $\lfloor x \rfloor$.
\end{enumerate}

Typically, tasks (\textit{what we want to achieve}) are about: given a function (\textit{relation between two sets}) we turn an input to an output. In some cases,
the input and output are not strings, so we need a way to represent them as a sequence of characters, following some encoding approaches.
\begin{example}
    e.g. Let's consider the following alphabet.
    $$S = \{a, b, c\}$$
    any single character of this alphabet can be encoded as:
    \begin{itemize}
        \renewcommand{\labelitemi}{}
        \item $a \rightarrow 00$
        \item $b \rightarrow 01$
        \item $c \rightarrow 10$
    \end{itemize}

    For instance, the original string $abbc$ becomes:
    $$00010110$$

    Note: we get an injective function but not a surjective function, since $11$ from the codomain corresponds to nothing in the domain.
\end{example}
\begin{definition}
    Given any natural number $n$, we get an encoding equal to transform it in a binary number:
    $$12 \rightarrow 1100$$
\end{definition}
\begin{definition}
    Given pairs of binary strings:
    $$\{0,1\}^* \times \{0,1\}^*$$
    where both $\{0,1\}^*$ can be a whatever sequence of binary numbers, we encode the resulting string $x \# y$ over the alphabet $\{0,1,\#\}$, where $\#$ acts as 
    \textbf{separator}.
\end{definition}
In computational theory, a task is interisting if it consists in computing a function from the alphabet $\{0,1\}^*$ to itself. An important function that starts
from $\{0,1\}^*$ and goes to $\{0,1\}^*$ are those whose range are strings of length exactly one, named \textbf{boolean functions}, also called \textbf{characteristic functions}.
\begin{definition}
    We identify the characteristic function such a function $f$ with a subset $L_f$ of $\{0,1\}^*$, defined as follows:
    $$L_f = \{x \in \{0,1\}^*|f(x) = 1\}$$
\end{definition}
\begin{definition}
    Any subset of $S^*$ (\textit{the set of any string of any dimension $n$}) is called \textbf{language}, indicated as $M$.
\end{definition}
These last definitions are extremely useful, because we can design a \textbf{decision problem} for a given language $M$ (\textit{subset of the alphabet $\{0,1\}^*$})
as task of computing a relation $f$ such that the language $M = L_f$. If we can derive such equivalance we get a resolution of the problem.  